\documentclass[11pt]{article}
\usepackage[margin=0.8in]{geometry}
\usepackage{amsmath,microtype}
\newcommand{\guideheading}[1]{\par\medskip\noindent{\large\bfseries #1}\par\smallskip}
\begin{document}
\begin{center}
{\Large Guide: Induction Steps for the Affine Pair}\par\medskip
Collatz proof project\quad---\quad September 5, 2026
\end{center}
\guideheading{The task}
The transfer problem asks whether reaching one passes from $3u+2$ to
$27u+20$ for every natural parameter $u$. The previous Lean result
shows that this would prove full Collatz. We now have two kinds of
verified certificates for particular parameter families.

\guideheading{A direct certificate}
For $u=36+64Q$, both partners reach exactly the same value after six
steps, for every natural $Q$. If the smaller partner reaches one,
the larger partner therefore does too. This is an unconditional proof
of the transfer implication on that progression.

\guideheading{A genuine induction step}
For the parameter family
\[
u=2308+4096Q,
\]
twelve steps transform the two partners into the same kind of pair at
\[
v=45+81Q<u.
\]
This reduces the question at $u$ to the question at a smaller parameter
$v$. Lean proves the endpoint identities, the strict inequality, and
the conditional transfer step for every quotient.

The smaller-parameter premise has not been removed. To turn these steps
into a complete induction, every parameter would need a direct proof,
a checked base case, or a reduction of this kind.

\guideheading{What the search found}
At depth 18, the exact search classifies 30,176 parameter residues by
direct merging and 1,267 by a smaller-parameter return. It leaves
230,701 unresolved. Each found rule describes all nonnegative quotients
on its binary arithmetic progression. The search stops at the first
certificate, so these counts describe that procedure rather than every
possible later certificate.

Independent tests replay all 3,236 saved certificates on several lifts
and compare small complete searches against direct orbits. The two
displayed infinite families and the general transport identities are
kernel checked; the full census is not a Lean completeness theorem.

\guideheading{Why ordinary convergence is not enough for a merge certificate}
At parameter zero, the partners are 2 and 20. Both reach one, but they
never have equal values at equal times. Lean checks this for all times.
Thus a search for aligned merging alone cannot settle every parameter.
Zero can be handled as a base case; the example does not disprove the
combined induction approach.

\guideheading{What remains open}
The missing step is coverage: a proof that every parameter can be
handled by a valid base case, direct merge, or contracting return.
The current families do not establish this. Collatz remains unproved.

\guideheading{Where to verify the work}
\begingroup\raggedright
Build \texttt{Collatz.AffinePairReturn} from \texttt{lean/}.
From the root, run \texttt{python3 analysis/affine\_pair\_search.py} or
\texttt{python3 -m unittest discover -s analysis -p test\_affine\_pair\_search.py}.
The technical paper states the exact lifting and induction conditions.
\par\endgroup
\end{document}
